% introduction_demonstration.tex
\subsection{The Present Demonstration}

What follows is one fully reproducible study exercising every exported function on the 50-item IPIP Big-Five Factor Markers, descendant of the lexical Big-Five program of \citet{goldberg-1990} and test bed of the work reviewed above \citep{casella-etal-2024,milano-etal-2025a,jung-seo-2025}. After the Method, the Results walk the workflow in analyst order: four embedding-based similarity encodings and a signed NLI matrix from the same 50 items, with the literature's reverse-keying problem treated as a modeling decision whose consequences are measured, not assumed; retention by multi-criterion consensus; a five-factor semantic model with Monte Carlo--calibrated diagnostics; interpretation via anchoring, projection, congruence, and jingle--jangle tools; and refinement via redundancy screening, a response-free short form, and vetting of newly written items. Finally, operationalizing the semantic--behavioral coherence relation proposed above, the semantic solution is compared at factor and item-pair level with a conventional exploratory factor analysis of more than 870{,}000 respondents to the same items \citep{openpsychometrics-2018}, showing, encoding by encoding, where the two relational structures converge and diverge and what the divergences reveal about the scale's language. Consistent with the convergence-with-divergence evidence above \citep{kojima-etal-2026,feraco-toffalini-2025}, the aim is calibration, not replacement: semantic factor analysis complements response-based analysis, hearing what a scale's language says before, and alongside, what its respondents say.
